% swim-times-stp.tex
% Software Test Plan for the swim-times program
% ISO/IEC/IEEE 29119-3:2021 Test Plan structure
% Compiled with: pdflatex swim-times-stp.tex

\documentclass[12pt]{article}

\usepackage{geometry}
\geometry{letterpaper, margin=1in}

\usepackage[T1]{fontenc}
\usepackage{hyperref}
\usepackage{booktabs}
\usepackage{array}
\usepackage{longtable}
\usepackage{enumitem}
\usepackage{titlesec}
\usepackage{fancyhdr}
\usepackage{xcolor}
\usepackage{listings}

\definecolor{castblue}{RGB}{0,48,135}
\definecolor{shade}{RGB}{245,246,248}

\hypersetup{colorlinks=true, linkcolor=castblue, urlcolor=castblue}

\pagestyle{fancy}
\fancyhf{}
\renewcommand{\headrulewidth}{0.4pt}
\fancyhead[L]{\small College Area Swim Team}
\fancyhead[R]{\small STP: swim-times}
\fancyfoot[C]{\thepage}

\titleformat{\section}{\large\bfseries\color{castblue}}{\thesection}{0.5em}{}[\titlerule]
\titleformat{\subsection}{\normalsize\bfseries}{\thesubsection}{0.5em}{}

\newcommand{\code}[1]{\texttt{#1}}
\newcommand{\req}[1]{\texttt{#1}}

\lstset{
  basicstyle=\ttfamily\small,
  backgroundcolor=\color{shade},
  frame=single,
  rulecolor=\color{gray!40},
  framesep=6pt,
  breaklines=true,
  columns=fullflexible,
  showstringspaces=false,
  xleftmargin=0pt
}

\begin{document}

% -----------------------------------------------------------------
\begin{titlepage}
  \centering
  \vspace*{2cm}
  {\Huge\bfseries\color{castblue} Software Test Plan\\[0.5em]}
  {\Large swim-times: Automated Swimmer-Times Retrieval System}\\[2em]
  {\large College Area Swim Team (CAST)}\\[0.5em]
  {\normalsize Prepared for: CAST Coaching Staff and Maintainers}\\[0.5em]
  {\normalsize Prepared by: CAST Technology Working Group}\\[2em]
  \begin{tabular}{ll}
    Document identifier: & CAST-STP-001                                       \\
    Revision:            & 6.0                                                \\
    Date:                & September 18, 2026                                 \\
    Status:              & Released                                           \\
    Supersedes:          & Revision 5.0, September 18, 2026                    \\
    Governing standard:  & ISO/IEC/IEEE 29119-3:2021 (Test Documentation)     \\
  \end{tabular}
  \vfill
  {\small\textit{This document is prepared in accordance with the principles
  of rational documentation articulated in D.\,L.\,Parnas and P.\,C.\,Clements,
  ``A Rational Design Process: How and Why to Fake It,''
  \textit{IEEE Transactions on Software Engineering}, vol.\,SE-12, no.\,2,
  pp.\,251--257, February 1986.}}
\end{titlepage}

\tableofcontents
\clearpage

% =================================================================
\section{Introduction}
% =================================================================

\subsection{Document identifier}

CAST-STP-001, revision 6.0.  This revision supersedes revision 5.0 in
its entirety.

\subsection{Purpose}

This plan defines the strategy, environment, design, and case-level
detail by which \code{swim-times} is verified against the requirements
of CAST-SRS-001.  It is the companion of the executable suites
\code{test-swim-times.sh} and \code{test-db.sh}: every case enumerated
in Section~\ref{sec:cases} is realised by one check in one of those
scripts, and every check in those scripts appears here.

\subsection{Rationale for this revision}\label{sec:why}

Revision 2.0 assumed something that is no longer true: that the
USA~Swimming data hub answers times queries from anonymous callers.  It
does not.  Since 2026 the hub serves reference data anonymously and
refuses every search and every times query with \code{403 Forbidden}
unless the caller is signed in.  The investigation is recorded in
CAST-DR-001.

That change breaks a test plan built on the old assumption in a
specific and damaging way: the suite reported eighteen failures whose
common cause was neither a defect nor a test error, and it reported
them in language that pointed at the program.  A suite that blames the
software for the world's changes is worse than no suite, because it
trains its reader to ignore it.

Revision 3.0 therefore adopts a different organising principle, stated
in full in Section~\ref{sec:tiers}: \textbf{the suite first determines
what it is permitted to observe, and then tests only what that
permits}, reporting every remaining case as \code{SKIP} with the actual
reason.  A \code{FAIL} in revision 3.0 always means a defect in
\code{swim-times}.

Revision 3.0 also added coverage that revision 2.0 lacked entirely: the
HTTP transport contract (\req{REQ-H-*}), the behaviour of the program
when it is refused (\req{REQ-E-*}), the correctness of the
cryptographic primitive behind the \code{rate-key} header
(\req{REQ-S-*}), and the self-diagnosis mode (\req{REQ-G-*}).  The
absence of the first of these is why the defect described in CAST-DR-001
was not caught by testing.

\subsection{Rationale for revision 4.0}

The program now signs itself in.  Revision 3.0 expected an operator to
extract a session from a browser by hand and export it, which meant the
live-data cases were skipped on every ordinary run --- twenty-one of
them.  Revision 4.0 of the software reads a user name and password from
a credentials file, completes the OpenID~Connect flow, and activates
the session; the live-data cases are therefore the \emph{normal} path,
and it is the refusal cases that skip.

The tier vocabulary changes accordingly (Section~\ref{sec:tiers}): the
question is no longer ``is this caller anonymous'' but ``did the
sign-in work, and if not, why''.  A new sub-process and a new
requirement family, \req{REQ-L-*}, cover the sign-in itself ---
including two properties that are security-relevant rather than merely
functional: that the credentials file is parsed and never sourced, and
that names outside the \code{USAS\_} prefix in it are ignored.

Revision 4.0 also records two defects the revised suite found on its
first authorised run, both introduced by revision 3.0 of the software:
\code{BestTimes} answers \code{404} for a stroke-and-distance pair the
swimmer has never swum, and the abort-on-failure logic added in
revision 3.0 treated that as fatal, ending the default all-events run
at the first such event (now \req{REQ-F-34}); and one CWEB chunk had
grown to twenty-six lines (\req{REQ-D-03}).

\subsection{Rationale for revision 5.0}

The program now uses one database --- \code{swimming} on
\code{100.77.243.69} --- named in a compiled-in constant, and adapts to
its normalised schema instead of writing a flat table of its own.  Two
things follow for testing.

First, there is no connection string to point at a scratch database, so
the suite tests against the real one.  It does this the only responsible
way: every write it makes is a clearly-marked fixture swimmer
(\code{ZZ swim-times \dots{} fixture}) removed again by an \code{EXIT}
trap, and the assertions about existing data are read-only counts.  The
suite adds and removes rows; it does not alter anyone else's.

Second, the database is shared with a sibling loader, so a passing run
is no longer evidence that this program wrote anything.  The cases are
framed accordingly: \req{REQ-F-73} asserts that a \emph{second} store
run adds nothing, which is the property that makes \code{-o store} safe
to schedule, and \req{REQ-F-76} asserts that the read joins through to
names, which is the property that would break first if the schema
moved.

\code{test-db.sh} is rewritten around the same principle: it verifies
the shape the program depends on --- the tables, the ten columns the
insert names, the absent default on \code{swim\_time\_id} that forces a
synthesised key, the natural-key constraint, all 31 event codes as
foreign-key targets, and that the \code{time\_standard} table rejects
the elite labels \code{db\_standard} maps to \code{NULL}.

\subsection{Rationale for revision 6.0}

Two defects were found by analysing the roster-to-database path, and
this revision adds the cases that pin them.  Both were found by
inspection rather than by a failing test, which is itself the finding:
the suite had no case that walked a roster containing an entry the
service could not resolve, and none that wrote a swimmer whose recorded
name differed in spelling from the one the service returns.

\req{REQ-F-77} and \req{REQ-F-77a} cover the first: a \code{404} from
the member search means ``no member matches this name'', and treating
it as a service failure aborted the whole roster walk at the first
stale entry, silently dropping every swimmer after it.

\req{REQ-F-78} through \req{REQ-F-80}, with \req{DB-19} and
\req{DB-20}, cover the second: the two writers of the shared database
learned their names from different feeds --- one recorded middle
initials, the other returns middle names --- so an exact-name match
found neither the row nor the conflict, and would have inserted a
duplicate person.  \req{DB-20} is the more important of the pair: it
asserts what the match must \emph{not} do.  Loosening a match is only
safe while it cannot cross into a row that has already been identified.

\subsection{Scope}

\paragraph{In scope.}  All requirements tagged \req{REQ-F-*},
\req{REQ-H-*}, \req{REQ-E-*}, \req{REQ-S-*}, \req{REQ-G-*},
\req{REQ-I-*}, \req{REQ-P-*}, \req{REQ-D-*}, and \req{REQ-A-*} in
CAST-SRS-001, plus the database-integrity cases \req{DB-01}
through \req{DB-10}.

\paragraph{Out of scope.}
\begin{itemize}
  \item The accuracy of the data USA~Swimming returns.  It is treated
        as authoritative.
  \item USA~Swimming's authorisation policy.  The suite observes and
        classifies it; it does not assert what it ought to be.
  \item The cwebmac macros and their typesetting choices, except that
        the woven document must typeset without \TeX{} errors
        (\req{REQ-A-04}).
\end{itemize}

\subsection{References}

\begin{enumerate}[leftmargin=*,label={[\arabic*]}]
\item CAST-SRS-001, \textit{Software Requirements Specification},
      \code{swim-times-srs.tex}.
\item CAST-SDD-001, \textit{Software Design Document},
      \code{swim-times-sdd.tex}.
\item CAST-CONOPS-001, \textit{Concept of Operations},
      \code{swim-times-conops.tex}.
\item CAST-DR-001, \textit{Defect Investigation and Corrective Action
      Report}, \code{swim-times-defect-report.tex}.
\item ISO/IEC/IEEE 29119-3:2021, \textit{Software and systems
      engineering --- Software testing --- Part 3: Test documentation}.
\item FIPS PUB 180-4, \textit{Secure Hash Standard}, NIST, 2015.
\item RFC 2104, \textit{HMAC: Keyed-Hashing for Message
      Authentication}, 1997; RFC 4231, \textit{Test Vectors for
      HMAC-SHA-256}, 2005.
\item D.\,L.\,Parnas and P.\,C.\,Clements, ``A Rational Design Process:
      How and Why to Fake It,'' \textit{IEEE Trans.\ Software
      Engineering}, SE-12(2), 1986.
\end{enumerate}

% =================================================================
\section{Test Context}
% =================================================================

\subsection{Test items}

\begin{itemize}
  \item \code{swim2/swim-times.w} --- the CWEB literate source, the
        canonical source of record.
  \item \code{swim2/swim-times.c} --- the tangled C source, a build
        artefact regenerated by the suite.
  \item \code{swim2/swim-times} --- the compiled executable.
  \item \code{swim2/swim-times.pdf} --- the woven documentation.
  \item \code{swim2/schema.sql} --- the Postgres mirror schema.
\end{itemize}

\subsection{Assumptions and constraints}

Revision 2.0 listed anonymous data-hub access as an assumption.  It is
now a \emph{variable}, observed at run time.  The remaining assumptions
are:

\begin{itemize}
  \item POSIX utilities (\code{awk}, \code{grep}, \code{sed},
        \code{cmp}, \code{curl}, and \code{otool} or \code{ldd}) are
        available.
  \item A C compiler and \code{ctangle}/\code{cweave} are available,
        natively or through the container wrapper \code{../cweb.sh}.
  \item \code{pdftex} is available for \req{REQ-A-04}.
  \item Where the database cases are to run: a reachable Postgres
        server holding a \code{swim-times} database.  Where it is not
        reachable, those cases \code{SKIP}.
  \item Where the live-data cases are to run: valid
        \code{USAS\_SUB\_ID} and \code{USAS\_SESSION\_ID} values.
        Where they are absent, those cases \code{SKIP} and the refusal
        cases run instead.
\end{itemize}

\subsection{Stakeholders}

\begin{center}
\begin{tabular}{ll}
  \toprule
  \textbf{Role}          & \textbf{Concern}                                \\
  \midrule
  Maintainer (developer) & Timely, correctly attributed regression signal  \\
  Coach (operator)       & Behavioural correctness of the CLI              \\
  Auditor                & Traceability from SRS to evidence               \\
  \bottomrule
\end{tabular}
\end{center}

% =================================================================
\section{Test Strategy}
% =================================================================

\subsection{The access-tier principle}\label{sec:tiers}

Before any behavioural case runs, the suite executes
\code{swim-times -o diag} and classifies the result into one of three
tiers.  The classification decides which cases are meaningful.

\begin{center}
\begin{tabular}{llp{5.6cm}}
  \toprule
  \textbf{Tier} & \textbf{Observed} & \textbf{Consequence} \\
  \midrule
  \textsc{authorized} & sign-in ok, queries 200
    & The normal outcome.  Live-data cases run; refusal cases
      \code{SKIP}. \\
  \textsc{nocreds} & no credentials found
    & Live-data cases \code{SKIP} naming the file to populate; refusal
      cases run. \\
  \textsc{authfail} & credentials present, sign-in failed
    & As above, with the reason ``check \code{USAS\_USER} and
      \code{USAS\_PASS}''. \\
  \textsc{locked} & sign-in ok, queries 401/403
    & The service refused an authenticated caller --- a change in its
      policy, or an activation failure.  Refusal cases run. \\
  \textsc{unreachable} & nothing answers
    & Both sets \code{SKIP}.  Static, CLI, crypto, sign-in-failure, and
      database cases still run. \\
  \bottomrule
\end{tabular}
\end{center}

\noindent Two properties follow, and they are the point of the design.
First, the suite gives a truthful verdict on any host, with or without
network, database, or account.  Second, a \code{FAIL} is unambiguous:
it is a defect in \code{swim-times}, never a statement about the
world.  \req{REQ-G-03} verifies that the classification itself works.

\subsection{Test sub-processes}

Nine sub-processes execute in order.  An earlier failure does not
suppress a later sub-process except where the artefact is missing
outright.

\begin{description}[style=nextline,leftmargin=1em]
  \item[1.\ Build.]  \code{make -B tangle compile} regenerates
        \code{swim-times.c} from the \code{.w} file and compiles it.  A
        failure here aborts the run: there is nothing to test.
  \item[2.\ Static and design constraints.]  Properties of the source
        and the binary --- warning-free compilation, module size,
        dependency set, absence of embedded credentials, HTTPS-only
        URLs, and a \TeX-error-free woven document.
  \item[3.\ HTTP transport contract.]  Source-level assertions that
        the transport layer reports status codes, distinguishes an
        empty body from a failure, sets timeouts, and sends no
        \code{Accept} header.  These are the regression guards for
        CAST-DR-001 root causes 2 and 3.
  \item[4.\ Cryptographic self-test.]  \code{-o selftest} against the
        FIPS 180-4 and RFC 4231 vectors, plus an independent
        cross-check of the expected digests.
  \item[5.\ Command-line interface.]  Exit codes, usage completeness,
        and mutually-exclusive option rejection.  Requires no network.
  \item[6.\ Access-tier determination.]  Section~\ref{sec:tiers}.
  \item[7.\ Sign-in.]  The credentials file, the sign-in failure paths,
        the activation step, and the absence of an anonymous fallback.
        Four of its seven cases construct their own fixture credentials
        and run in any tier.
  \item[7b.\ Refusal handling.]  Runs when the service declines the
        caller: the program's behaviour in that event.
  \item[8.\ Live data retrieval.]  Runs in tier \textsc{authorized},
        which is now the expected outcome of an ordinary run.
  \item[9.\ Database persistence.]  Schema conformance, the duplicate
        guard, and the offline read path, against seeded fixtures.
\end{description}

\subsection{Test design techniques}

\begin{description}[style=nextline,leftmargin=1em]
  \item[Specification-based (black box).]  CLI behaviour, exit codes,
        output format, and retrieved content are exercised through the
        published interface only.
  \item[Structure-based (white box), by inspection.]  The
        \req{REQ-H-*} and \req{REQ-D-*} cases assert properties of the
        literate source that cannot be observed from outside --- that
        the status is read, that an empty body is normalised, that no
        chunk exceeds twenty-four lines.  These are regression guards:
        their purpose is to fail loudly if a future edit reintroduces a
        defect whose external symptom is subtle.
  \item[Known-answer testing.]  SHA-256 and HMAC-SHA256 are verified
        against published vectors rather than against themselves.
  \item[Fixture-based testing.]  Database cases seed deterministic rows
        and remove them afterwards, so they neither depend on what the
        live service happens to return nor leave residue behind.
  \item[Error guessing.]  The refusal cases were derived from the
        specific failure modes observed in CAST-DR-001, including the
        one that produced 638 doomed requests per run.
  \item[Negative testing with a constructed fixture.]  \req{REQ-L-03}
        and \req{REQ-L-04} write a credentials file containing a
        command substitution and a non-\code{USAS\_} assignment, then
        assert that the substitution did not run and the assignment was
        not adopted.  Both properties are invisible from the program's
        normal output and can only be established this way.
\end{description}

\subsection{Test environment}

\begin{center}
\begin{tabular}{ll}
  \toprule
  \textbf{Element}        & \textbf{Requirement}                        \\
  \midrule
  Operating system        & POSIX (macOS or Linux)                      \\
  Compiler                & C99, \code{-O2 -Wall}                       \\
  Libraries               & libcurl, libpq                              \\
  CWEB tools              & \code{ctangle}, \code{cweave} (native or container) \\
  Typesetting             & \code{pdftex}, \code{pdflatex}              \\
  Database (optional)     & \code{swimming} on \code{100.77.243.69},
                            reachable via libpq                        \\
  Network (optional)      & Outbound HTTPS to \code{times-api.usaswimming.org} \\
  Credentials (optional)  & \code{USAS\_USER} and \code{USAS\_PASS} in
                            \code{\$HOME/.usas-env}, or a session in
                            \code{USAS\_SUB\_ID}/\code{USAS\_SESSION\_ID} \\
  \bottomrule
\end{tabular}
\end{center}

\noindent The three optional elements are what the tiering exists to
handle; none is a precondition for a useful run.

\subsection{Test data}

\begin{itemize}
  \item \textbf{Roster data} is the \code{SWIMMERS} table compiled into
        the program.
  \item \textbf{The probe member} is Katie Ledecky's public
        \code{memberId} \code{6CD35348E5824C}, used by \code{-o diag}
        and \req{REQ-F-18}.  It is a matter of public record and is
        read-only.
  \item \textbf{Database fixtures} are a swimmer named
        \code{ZZ swim-times \dots{} fixture (delete me)} and the swims
        attached to it.  The name cannot collide with a real swimmer,
        and \code{swim} cascades on delete, so removing the swimmer
        removes everything the fixture created.  Both suites install an
        \code{EXIT} trap so the fixture is removed even if a case fails
        or the run is interrupted.
  \item \textbf{Credentials fixtures} are a temporary
        \code{USAS\_ENV\_FILE} written into the suite's scratch
        directory, carrying deliberately invalid credentials, a command
        substitution used as a tripwire, and a bogus
        \code{SWIM\_TIMES\_PGCONNINFO}.  No real credential is ever
        written by the suite.
  \item \textbf{Cryptographic vectors} are the published values for
        SHA-256 of the empty string, of \code{"abc"}, and of 199
        \code{'a'} bytes, and HMAC-SHA256 under the empty key and under
        \code{"key"}.
\end{itemize}

\subsection{Item pass/fail criteria}

\begin{itemize}
  \item A case \textbf{passes} when its assertion holds.
  \item A case \textbf{fails} when its assertion does not hold.  Every
        failure is a defect in \code{swim-times} or in its
        documentation.
  \item A case is \textbf{skipped} when a stated precondition is not
        met on the test host.  The reason is printed with the skip.  A
        skip is not a failure and does not affect exit status.
  \item The suite \textbf{passes} when no case fails.  Its exit status
        is 0 in that event and 1 otherwise.
\end{itemize}

\noindent A release is blocked by any failure.  A release is
\emph{not} blocked by skips, but a release made while the live-data
cases are skipped carries the caveat that live retrieval was not
demonstrated on this run --- which the tester shall record.

\subsection{Suspension and resumption}

Testing is suspended if the build sub-process fails, since every later
case depends on the executable.  It resumes when a build is obtained.
No other condition suspends the run: the tiering absorbs a missing
network, a missing database, and missing credentials.

\subsection{Test deliverables}

\begin{itemize}
  \item This plan (CAST-STP-001).
  \item \code{test-swim-times.sh} --- the requirements suite.
  \item \code{test-db.sh} --- the database-integrity suite.
  \item The console transcript of each run, which is the test log.
  \item CAST-DR-001, where a run's findings warrant an investigation.
\end{itemize}

% =================================================================
\section{Test Cases and Traceability}\label{sec:cases}
% =================================================================

\subsection{Static and design constraints}

\begin{center}
\begin{longtable}{p{2.1cm}p{5.2cm}p{6.2cm}}
  \toprule
  \textbf{ID} & \textbf{Requirement} & \textbf{Method} \\
  \midrule\endhead
  \req{BUILD}     & Tangle and compile succeed
    & \code{make -B tangle compile}; executable exists. \\
  \req{REQ-D-01}  & Warning-free build
    & No \code{warning:} in the build log under \code{-O2 -Wall}. \\
  \req{REQ-D-02}  & \code{.c} is a regenerable artefact
    & \code{swim-times.c} is newer than \code{swim-times.w} after the
      build. \\
  \req{REQ-D-03}  & No CWEB chunk exceeds 24 lines
    & \code{awk} audit over the \code{.w} file. \\
  \req{REQ-D-04}  & Dependency set is \{libc, libcurl, libpq, TLS\}
    & \code{otool -L} / \code{ldd} against an allowlist.  In
      particular no crypto library may appear: SHA-256 is in-tree. \\
  \req{REQ-D-11}  & No embedded credential
    & Pattern scan for bearer tokens and JWTs in the source. \\
  \req{REQ-I-20}  & HTTPS-only service URLs
    & \code{https://} base URLs present, no \code{http://} hub URL. \\
  \req{REQ-A-04}  & Documentation typesets cleanly
    & \code{cweave} then \code{pdftex}; PDF exists and the log
      contains zero \code{!} errors. \\
  \bottomrule
\end{longtable}
\end{center}

\subsection{HTTP transport contract (regression guards, CAST-DR-001)}

\begin{center}
\begin{longtable}{p{2.1cm}p{5.2cm}p{6.2cm}}
  \toprule
  \textbf{ID} & \textbf{Requirement} & \textbf{Method} \\
  \midrule\endhead
  \req{REQ-H-01} & \code{http\_request} exposes the HTTP status
    & Signature carries a \code{long *status} out-parameter. \\
  \req{REQ-H-02} & The status comes from libcurl
    & \code{CURLINFO\_RESPONSE\_CODE} appears in the source. \\
  \req{REQ-H-03} & An empty body is not a failure
    & The normalisation \code{calloc(1,\,1)} is present, so
      \code{NULL} means ``no response'' and nothing else. \\
  \req{REQ-H-04} & Timeouts are configured
    & \code{CURLOPT\_TIMEOUT} and \code{CURLOPT\_CONNECTTIMEOUT}
      present. \\
  \req{REQ-H-05} & No \code{Accept} header is sent
    & \code{"Accept: application/json"} absent from the source.
      Sending it makes the service double-encode its payload. \\
  \req{REQ-H-06} & Every response site checks the status
    & Count of \code{status != 200} guards $\geq$ the number of
      response-consuming call sites. \\
  \bottomrule
\end{longtable}
\end{center}

\subsection{Cryptographic self-test}

\begin{center}
\begin{longtable}{p{2.1cm}p{5.2cm}p{6.2cm}}
  \toprule
  \textbf{ID} & \textbf{Requirement} & \textbf{Method} \\
  \midrule\endhead
  \req{REQ-S-01} & Known-answer vectors pass
    & \code{-o selftest} exits 0 and reports ``all vectors passed''. \\
  \req{REQ-S-02} & At least five vectors are exercised
    & Count of \code{ok} verdicts in the self-test output. \\
  \req{REQ-S-03} & Expected digests are independently derived
    & The HMAC vector in the source matches Python \code{hashlib}.
      Skipped where \code{python3} is unavailable. \\
  \bottomrule
\end{longtable}
\end{center}

\subsection{Command-line interface}

\begin{center}
\begin{longtable}{p{2.1cm}p{5.2cm}p{6.2cm}}
  \toprule
  \textbf{ID} & \textbf{Requirement} & \textbf{Method} \\
  \midrule\endhead
  \req{REQ-F-03} & Empty invocation prints usage, exits 2
    & Run with no arguments; check status and output. \\
  \req{REQ-F-04} & Unknown flag prints usage, exits 2
    & Run with \code{-Z}. \\
  \req{REQ-F-05} & Usage documents every keyword
    & \code{fastest}, \code{csv}, \code{store}, \code{offline},
      \code{diag}, \code{selftest}, \code{-m}, and both credential
      variables appear in the usage text. \\
  \req{REQ-F-06} & Usage lists the whole roster
    & Declared count $\geq 29$, generated from \code{SWIMMERS}. \\
  \req{REQ-F-07} & \code{diag} and \code{offline} are exclusive
    & Combination exits 2 with an explanatory message. \\
  \bottomrule
\end{longtable}
\end{center}

\subsection{Sign-in}

\begin{center}
\begin{longtable}{p{2.1cm}p{5.2cm}p{6.2cm}}
  \toprule
  \textbf{ID} & \textbf{Requirement} & \textbf{Method} \\
  \midrule\endhead
  \req{REQ-L-01} & Absent credentials are named precisely
    & Run with \code{USAS\_ENV\_FILE} pointing at a non-existent path
      and the credential variables unset; expect exit 77 and a message
      naming both the file and the variables. \\
  \req{REQ-L-02} & Rejected credentials are reported in plain language
    & Run with a deliberately wrong password; expect exit 77 and
      ``rejected the credentials'', \emph{not} the \code{HTTP 200} the
      identity provider actually answers with. \\
  \req{REQ-L-03a} & Credentials are read from the file
    & A fixture file supplies a distinctive user name; the failure
      message must quote it. \\
  \req{REQ-L-03} & The file is parsed, not sourced
    & The same fixture contains
      \code{USAS\_SIDE\_EFFECT="\$(touch \dots)"}; the marker file must
      not exist afterwards. \\
  \req{REQ-L-04} & Non-\code{USAS\_} assignments are ignored
    & The same fixture sets a bogus
      \code{SWIM\_TIMES\_PGCONNINFO}; an offline read with the real
      variable unset must still succeed, proving the built-in default
      was used. \\
  \req{REQ-L-05} & \code{offline} needs no credentials
    & Offline read with every credential variable unset and no file. \\
  \req{REQ-L-06} & The session is activated
    & \code{GetDataHubSecurityInfoForIdp} and \code{activate\_session}
      appear in the source.  Its behavioural counterpart is
      \req{REQ-G-03} reaching tier \textsc{authorized}: without the
      activation the probes return \code{401}. \\
  \req{REQ-L-07} & No anonymous data path
    & Every non-\code{offline} run is gated on \code{ensure\_session}. \\
  \bottomrule
\end{longtable}
\end{center}

\subsection{Access-tier determination}

\begin{center}
\begin{longtable}{p{2.1cm}p{5.2cm}p{6.2cm}}
  \toprule
  \textbf{ID} & \textbf{Requirement} & \textbf{Method} \\
  \midrule\endhead
  \req{REQ-G-01} & \code{diag} probes every access class
    & At least five probe lines; exit status 0. \\
  \req{REQ-G-02} & \code{diag} reports the identity presented
    & \code{Usas-Sub-Id} and \code{Device-Id} appear in the output. \\
  \req{REQ-G-03} & The five tiers are distinguishable
    & The suite classifies the run from \code{diag} output alone and
      reports which tier it selected. \\
  \bottomrule
\end{longtable}
\end{center}

\subsection{Refusal handling}

\begin{center}
\begin{longtable}{p{2.1cm}p{5.2cm}p{6.2cm}}
  \toprule
  \textbf{ID} & \textbf{Requirement} & \textbf{Method} \\
  \midrule\endhead
  \req{REQ-E-01} & A refusal exits \code{RC\_NOPERM} (77)
    & Single-swimmer run in a non-authorized tier; check status. \\
  \req{REQ-E-02} & The diagnostic is specific
    & \code{stderr} names \code{HTTP 403 Forbidden} and the endpoint,
      and does not contain the superseded text
      ``person lookup request failed''. \\
  \req{REQ-E-03} & The remedy is stated once per run
    & \code{USAS\_SESSION\_ID} appears exactly once in \code{stderr}. \\
  \req{REQ-E-04} & A refusal aborts promptly
    & Full-roster run emits at most two \code{Error:} lines rather
      than replaying the refusal across 29 swimmers. \\
  \req{REQ-F-18} & The known-\code{memberId} short cut works
    & \code{-m 6CD35348E5824C} resolves the member through
      \code{GetMember} without a search.  Runs in every tier with a
      network, since \code{GetMember} is open to anonymous callers. \\
  \bottomrule
\end{longtable}
\end{center}

\subsection{Live data retrieval (tier \textsc{authorized})}

\begin{center}
\begin{longtable}{p{2.1cm}p{5.2cm}p{6.2cm}}
  \toprule
  \textbf{ID} & \textbf{Requirement} & \textbf{Method} \\
  \midrule\endhead
  \req{REQ-F-10} & Swimmer selection restricts output
    & Only the selected swimmer appears. \\
  \req{REQ-F-11} & Multiple ids select multiple swimmers
    & Count of \code{Swimmer:} lines $\geq 2$. \\
  \req{REQ-F-12} & \code{stella} resolves correctly
    & Resolved name contains \code{Evans}. \\
  \req{REQ-F-13} & \code{ledecky10} resolves correctly
    & Resolved name contains \code{Ledecky}. \\
  \req{REQ-F-14/15/16} & Age windows are declared
    & The three Ledecky windows appear on their \code{Swimmer}
      entries.  Checked statically, so these run in every tier. \\
  \req{REQ-F-16a} & An age-windowed row carries a date
    & Row matches the ISO date pattern. \\
  \req{REQ-F-20} & A single \code{-e} restricts to one event
    & Exactly one event section. \\
  \req{REQ-F-21} & All 31 events by default
    & Thirty-one event sections when \code{-e} is omitted. \\
  \req{REQ-F-22} & The event list is honoured
    & Section count equals the request. \\
  \req{REQ-F-30} & \code{memberId} resolution succeeds
    & A \code{Swimmer:} line and an event section are produced. \\
  \req{REQ-F-31} & The times query returns results
    & At least one event section with data. \\
  \req{REQ-F-32} & Rows carry time, date, and meet
    & A row matches the full-row pattern. \\
  \req{REQ-F-33} & Per-pair caching is present
    & The skip-if-already-fetched chunk exists. \\
  \req{REQ-F-34} & A \code{404} means ``no such time'', not failure
    & The all-events run for a young swimmer must report all
      thirty-one sections, most of them \code{(no times found)}.  This
      is \req{REQ-F-21}: before the fix it saw zero. \\
  \req{REQ-F-40} & Rows are fastest-first
    & The emitted order equals the sorted order. \\
  \req{REQ-F-41} & \code{fastest} prints one row
    & At most one data row. \\
  \req{REQ-F-42} & Dates are ISO \code{YYYY-MM-DD}
    & Pattern match on every date field. \\
  \req{REQ-F-50} & Table mode prints heading and header
    & Column header present. \\
  \req{REQ-F-51} & CSV header is the six-column contract
    & Exact string comparison. \\
  \req{REQ-F-52} & Every CSV row carries the swimmer name
    & No data row lacks a leading quoted name field. \\
  \req{REQ-P-01} & Full refresh under five minutes
    & Extrapolated from the measured single-query time. \\
  \req{REQ-P-02} & Single query under ten seconds
    & Wall-clock measurement. \\
  \req{REQ-A-01} & Repeated runs are identical
    & \code{cmp} of two consecutive runs. \\
  \bottomrule
\end{longtable}
\end{center}

\subsection{The shared database}

Run by \code{test-swim-times.sh} (\req{REQ-F-7x}) and in greater depth
by \code{test-db.sh} (\req{DB-*}).  Both use the compiled-in connection
string, because the program admits no other.

\begin{center}
\begin{longtable}{p{2.1cm}p{5.2cm}p{6.2cm}}
  \toprule
  \textbf{ID} & \textbf{Requirement} & \textbf{Method} \\
  \midrule\endhead
  \req{REQ-F-70} & \code{swim} carries every column the insert names
    & Catalogue query for the ten columns of
      \code{db\_insert\_row}. \\
  \req{REQ-F-71} & The offline read joins through to names
    & \code{-o offline,csv} returns rows whose swimmer and meet fields
      are populated, proving the \code{swim}$\to$\code{swimmer} join
      and the \code{meet} outer join. \\
  \req{REQ-F-72} & Offline mode performs no network I/O
    & curl initialisation is guarded by the offline flag. \\
  \req{REQ-F-73} & Writes are idempotent
    & Two consecutive \code{-o store} runs; the second must add no
      rows.  This is what makes \code{-o store} safe to schedule
      against a database somebody else also writes. \\
  \req{REQ-F-75} & A run that cannot sign in writes nothing
    & Store run with the credentials removed; row count unchanged and
      exit 77. \\
  \req{REQ-F-77} & An absent swimmer is skipped, not fatal
    & \code{-m} with an unknown member id: expect a warning, the
      end-of-run tally, and exit 69. \\
  \req{REQ-F-77a} & The walk continues past one
    & Structural: the dispatcher must \code{continue} rather than
      \code{break}.  A roster entry that fails to resolve cannot be
      injected from outside the binary, so this is inspected. \\
  \req{REQ-F-78} & The loose match is confined to unlinked rows
    & Structural: both fuzzy tiers carry \code{member\_id IS NULL}. \\
  \req{REQ-F-76} & The surrogate swim id is stable
    & \code{swim\_surrogate\_id} hashes the natural key, and the insert
      uses \code{ON CONFLICT DO NOTHING}.  An unstable id would insert
      the same swim again under a fresh primary key on every run. \\
  \midrule
  \req{DB-01} & The database is reachable
    & \code{SELECT 1} against
      \code{swimming@100.77.243.69}. \\
  \req{DB-02} & The program names that same database
    & The connection string in \code{swim-times.w} matches the one
      under test.  A mismatch would make every case below a report
      about the wrong server. \\
  \req{DB-03} & It cannot be redirected
    & No \code{getenv} for \code{SWIM\_TIMES\_PGCONNINFO} remains.  The
      name may still appear in the prose explaining its withdrawal. \\
  \req{DB-04} & The five tables the program names exist
    & \code{swimmer}, \code{meet}, \code{swim}, \code{event},
      \code{time\_standard}. \\
  \req{DB-05/06/07} & Their columns exist
    & Per-table catalogue queries for the columns the insert and the
      swimmer/meet resolution name. \\
  \req{DB-08} & \code{swim\_time\_id} has no default
    & Which is \emph{why} the program synthesises one.  Were a default
      to appear, \code{swim\_surrogate\_id} would become unnecessary
      and the code should be revisited --- so this is a notification,
      not merely a guard. \\
  \req{DB-09} & The natural-key constraint is present
    & Matched by column set: \code{(swimmer\_key, event\_code,
      swim\_time, swim\_date, meet\_key)}. \\
  \req{DB-10} & All 31 event codes are valid FK targets
    & Every code the program can request exists in \code{event}; one
      absent would make that event's rows unwritable. \\
  \req{DB-11} & \code{time\_standard} holds the seven mapped levels
      and no elite labels
    & Confirms both halves of \code{db\_standard}'s contract. \\
  \req{DB-12} & Repeated identical inserts collapse
    & Three inserts of one fixture swim yield one row. \\
  \req{DB-13} & A distinct swim is not merged
    & A second, different fixture swim yields a second row. \\
  \req{DB-14} & An elite label is rejected by the FK
    & Inserting \code{'Summer Jrs'} must fail where the same row with
      a \code{NULL} standard succeeds --- attributing the failure to the
      label rather than to the rest of the row. \\
  \req{DB-19} & The loose match claims an unlinked row
    & A fixture swimmer with a middle initial and no member id is
      matched by a first/last query spelling the middle name in full ---
      the case that produced duplicates. \\
  \req{DB-20} & \dots{} and never a linked one
    & The same query against the same fixture, once it carries a
      \emph{different} member id, must match nothing.  This is the
      guard that makes the fuzziness safe. \\
  \req{DB-15/16/17} & The read path works end to end
    & \code{-o offline,csv} exits 0 with empty \code{stderr}, emits the
      six-column header, and returns rows with no empty swimmer
      name. \\
  \req{DB-18} & Contents are reportable
    & Swimmer, meet, and swim counts, broken down by
      \code{source\_endpoint} and \code{auth\_mode} so the two loaders'
      contributions are distinguishable. \\
  \bottomrule
\end{longtable}
\end{center}

% =================================================================
\section{Execution and Reporting}
% =================================================================

\subsection{Invocation}

\begin{lstlisting}
cd swim2
make test                      # or: ./test-swim-times.sh
make db-test                   # or: ./test-db.sh

./test-swim-times.sh --no-network   # static, CLI, crypto, DB only
./test-swim-times.sh --no-db        # skip the database sub-process
\end{lstlisting}

\noindent No credential handling is needed at the command line: the
program signs itself in from \code{\$HOME/.usas-env}, so an ordinary
run reaches tier \textsc{authorized} on its own.  To test a different
account, point \code{USAS\_ENV\_FILE} at another file; to test the
unauthenticated paths deliberately, unset the credential variables and
point it at a non-existent path.

\subsection{Result format}

Each case emits one line:

\begin{lstlisting}
PASS REQ-H-03     empty response body normalised to "" rather than NULL
FAIL REQ-E-01     refused run exits 77 (RC_NOPERM), got 1
SKIP REQ-F-31     data hub requires sign-in; set USAS_SUB_ID and USAS_SESSION_ID
\end{lstlisting}

\noindent followed by a summary count.  Exit status is 0 when no case
failed.  The transcript is the test log; no external tool is required
to interpret it.

\subsection{Reference result}

The baseline recorded on release of this revision, on a host with
network, database, and a working USA~Swimming account (tier
\textsc{authorized}):

\begin{center}
\begin{tabular}{lrrr}
  \toprule
  \textbf{Suite} & \textbf{Pass} & \textbf{Fail} & \textbf{Skip} \\
  \midrule
  \code{test-swim-times.sh} & 67 & 0 & 4 \\
  \code{test-db.sh}         & 20 & 0 & 0 \\
  \bottomrule
\end{tabular}
\end{center}

\noindent The four skips are the refusal cases \req{REQ-E-01} through
\req{REQ-E-04}, which cannot run while the service is answering.  On
the same host with \code{--no-network} the suite records 31 passes and
32 skips, and still no failures.

For comparison, revision 3.0 of this plan recorded 38 passes and 21
skips on the same host --- the difference is the twenty-one live-data
cases that signing in makes reachable, plus the \req{REQ-L-*} and
\req{REQ-F-76} cases added since.  \code{test-db.sh} grew from 10 cases
to 18 when it was rewritten around the shared schema.

% =================================================================
\section{Risk Register}
% =================================================================

\begin{center}
\begin{longtable}{p{4.0cm}p{1.5cm}p{1.5cm}p{6cm}}
  \toprule
  \textbf{Risk} & \textbf{Likeli-} & \textbf{Impact} & \textbf{Mitigation} \\
                & \textbf{hood}    &                 & \\
  \midrule\endhead
  Further narrowing of the data-hub interface
      & High & High & \req{REQ-G-01..04} detect and classify it on
                      every run; the Postgres mirror is treated as the
                      system of record. \\
  A new status is mistaken for an error
      & Med  & Med  & \req{REQ-F-34} and \req{REQ-F-77} pin the two
                      already found, both \code{404}; \req{REQ-F-21}
                      and \req{REQ-F-77a} are the end-to-end checks
                      that would catch another. \\
  A roster query goes stale
      & High & Low  & \req{REQ-F-77} keeps it to one skipped swimmer,
                      named in the tally.  The remedy is to give the
                      entry a \code{member\_id}, as \code{kenny} and
                      \code{keith} now carry. \\
  Two swimmers merged by a loose name match
      & Low  & High & \req{DB-20} asserts the guard directly; the
                      fuzzy tiers cannot touch a row that already has
                      a member id. \\
  \code{GetMember} also closes
      & Med  & Med  & \req{REQ-F-18} fails immediately, distinguishing
                      this from the broader lockout. \\
  Session expires mid-run
      & Low  & Low  & Each run signs in afresh, so the window is one
                      run; a failure gives partial output plus
                      \code{RC\_NOPERM}. \\
  Sign-in flow changes shape
      & Med  & High & \req{REQ-N-04} is implemented as five discrete
                      steps, each reporting its own failure, so a
                      change names the step that broke.
                      \req{REQ-L-01,02} exercise the failure paths. \\
  Activation coupling removed or changed
      & Med  & Med  & \req{REQ-L-06} is structural and \req{REQ-G-03}
                      behavioural; a change surfaces as tier
                      \textsc{locked} with \code{401} probes. \\
  Stored password becomes stale
      & Med  & Low  & \req{REQ-L-02}'s message is the one the operator
                      will see; it names the two variables to fix. \\
  \code{Device-Id} format changes
      & Low  & High & The service answers \code{400 Invalid Device-Id
                      format}, which \req{REQ-G-01} surfaces as a
                      non-200 probe. \\
  Response schema change
      & Low  & High & \req{REQ-F-32} detects parse failure on the
                      first authorized run. \\
  A future edit reintroduces a fixed defect
      & Med  & High & \req{REQ-H-01..06} are structural guards written
                      specifically against CAST-DR-001. \\
  Test host lacks libcurl or libpq
      & Low  & High & The build sub-process fails fast and testing is
                      suspended. \\
  The shared database's schema moves
      & Med  & High & \req{DB-04} through \req{DB-11} check the shape
                      the program depends on, column by column, and
                      name what changed. \\
  A fixture is left behind after a failed run
      & Low  & Low  & Both suites remove the fixture from an
                      \code{EXIT} trap, so an interrupt or a failing
                      case still cleans up. \\
  The sibling loader and this program disagree
      & Med  & Med  & \code{source\_endpoint} and \code{auth\_mode}
                      distinguish their rows (\req{DB-18}); neither
                      deletes or updates the other's. \\
  \bottomrule
\end{longtable}
\end{center}

% =================================================================
\section{Approvals}
% =================================================================

\begin{center}
\begin{tabular}{p{5cm}p{5cm}p{3cm}}
  \toprule
  \textbf{Role} & \textbf{Name} & \textbf{Date} \\
  \midrule
  Prepared by (Test Lead)   & \rule{4.5cm}{0.4pt} & \rule{2.5cm}{0.4pt} \\[1.2em]
  Reviewed by (Maintainer)  & \rule{4.5cm}{0.4pt} & \rule{2.5cm}{0.4pt} \\[1.2em]
  Approved by (CAST)        & \rule{4.5cm}{0.4pt} & \rule{2.5cm}{0.4pt} \\
  \bottomrule
\end{tabular}
\end{center}

\end{document}
